% Generated from locale/tr/content/set-theory/replacement/extrinsic.tex; do not hand-edit.


\olfileid[tr]{sth}{replacement}{extrinsic}
\olsection[Dışsal Değerlendirmeler]{Yerine Koyma Hakkında Dışsal Değerlendirmeler}

Yerine Koyma'yı gerekçelendirmeye yönelik \emph{dışsal} bir girişimi ele
alarak başlayalım. Boolos şöyle bir girişim önerir:
\begin{quote}
  [\ldots] yerine koyma aksiyomlarını benimsemenin nedeni oldukça basittir:
  arzu edilir birçok sonuçları ve (görünüşe göre) arzu edilmez hiçbir
  sonuçları yoktur. Yinelemeli anlayışa ilişkin teoremlerin yanı sıra bu
  sonuçlar, ideal olmasa da doyurucu bir sonsuz sayılar kuramını ve iyi
  temellenmiş bağıntılar üzerindeki tümevarımlı tanımları gerekçelendiren son
  derece arzu edilir bir sonucu içerir.
  \citep[229]{Boolos1971}
\end{quote}		
Boolos'un düşüncesinin özü, Yerine Koyma'yı meyveleriyle gerekçelendirmemiz
gerektiğidir. Sözünü ettiği belirli meyveler de son birkaç bölümde
tartıştığımız şeylerdir. Yerine Koyma, von Neumann sıra sayılarının bir iyi
sıralama tipi fikri için mükemmel temsilciler olduğunu ispatlamamızı sağladı
(bu, ``ideal olmasa da doyurucu sonsuz sayılar kuramımızdır''). Yerine Koyma
ayrıca $V_\alpha$'ları tanımlamamızı, rank kavramını kurmamızı ve
$\in$-Tümevarım'ı ispatlamamızı sağladı (bunlar ``yinelemeli anlayışa ilişkin
teoremlerimiz''dir). Son olarak Yerine Koyma, Sonlu Ötesi Özyineleme
Teoremi'ni ispatlamamızı sağlar (bu da ``iyi temellenmiş bağıntılar üzerindeki
tümevarımlı tanımlar''dır).

Bunlar gerçekten arzu edilir sonuçlardır. Fakat bu sonuçlar Yerine Koyma'yı
\emph{gerekçelendirmeye} yeter mi? \emph{Hayır}. Ya da en azından doğrudan
yetmez.

Basit bir sorun şöyledir. Yerine Koyma'nın bazı arzu edilir sonuçlarını
belirttik; ancak bunların çoğunu başka yollarla elde edebilirdik. Bu durum
olması gerektiği kadar iyi bilinmediğinden açıklamak üzere duraklamalıyız.

Düzey Teorisi denen, kısaca $\LT$ diyeceğimiz basit bir kümeler teorisi
vardır.\footnote{$\LT$'nin ilk sürümleri \citet{Montague1965} ve
\citet{Scott1974} tarafından sunulmuştur; \citet{Potter2004} bunu
sadeleştirip kitap uzunluğunda ele almış; yakın zamanda \citet{ButtonLT1}
$\LT$'yi daha da sadeleştirmiştir.} $\LT$'nin aksiyomları yalnızca
Kaplamsallık, Ayırma ve her kümenin bir \emph{düzeyin} alt kümesi olduğu
iddiasıdır; burada ``düzey'', düzeyler dostlarımız $V_\alpha$'lar gibi
davranacak biçimde ustaca tanımlanır. Dolayısıyla $\ZF$, $\LT$'yi ispatlar;
ama $\LT$, $\ZF$'den \emph{çok} daha zayıftır. Aslında $\LT$ bize Çiftleri,
Kuvvet Kümelerini, Sonsuzluğu ya da Yerine Koyma'yı vermez. $\LT$'ye
Sonsuzluk ile Kuvvet Kümelerini ekleyerek elde edilen teoriye $\Zr$ diyelim;
bu teori Çiftleri de verir, dolayısıyla $\Zr$ en az $\Z$ kadar güçlüdür.
Fakat aslında $\Zr$, her kümenin bir rankı olduğu iddiasını da eklediği için
$\Z$'den kesin olarak daha güçlüdür (ona $\Zr$ demeyi önermemin nedeni de
budur). Nitekim $\Zr$ şunları sağlar: tamamen doyurucu bir sıra sayıları
teorisi; hiyerarşiyi iyi sıralanmış aşamalara katmanlaştıran sonuçlar;
$\in$-Tümevarım'ın bir ispatı ve Sonlu Ötesi Özyineleme'nin bir
\emph{sürümü}.

Kısacası Boolos bunu bilmese de sözünü ettiği bütün arzu edilir sonuçlara
Yerine Koyma \emph{olmadan} ulaşılabilirdi; yalnızca $\Z$ yerine $\Zr$
kullanması gerekirdi.

(Bütün bunlar doğruysa neden geleneksel yolu izleyip size $\LT$ ile $\Zr$
yerine $\ZF$ öğrettik? Bunun iki nedeni vardır. Birincisi, sırf tarihsel
nedenlerle $\LT$ ile başlamak oldukça standart dışıdır; sizi küme teorisinin
daha standart tartışmalarını okuyabilecek biçimde donatmak istedik. İkincisi,
$\LT$ ile $\Zr$'yi değerlendirmeye hazır olduğunuzda doğrudan
\citealt{Potter2004} ve \citealt{ButtonLT1} okuyabilirsiniz.)

Elbette $\Zr$, $\ZF$'den kesin olarak daha zayıf olduğundan $\ZF$'nin
ispatladığı ama $\Zr$'nin açık bıraktığı sonuçlar vardır. Dolayısıyla bu
sonuçlardan birine işaret ederek Yerine Koyma'yı dışsal gerekçelerle
gerekçelendirmeye çalışılabilir. Ancak $\Zr$'yi nasıl kullanacağınızı
bildiğinizde, hem (a) başka türlü değil de Yerine Koyma tarafından karara
bağlanan hem de (b) sezgisel olarak doğru olan şeylere ilişkin çok sayıda
örnek bulmak oldukça güçtür. (Bu konuda daha fazlası için
\citealt[\S13.2]{Potter2004} kaynağına bakın.)

Sonuç şudur. Yerine Koyma için ikna edici bir dışsal gerekçelendirme sunmak
üzere Yerine Koyma olmadan \emph{elde edilemeyecek} bir sonuç bulmak
gerekir. Bu da kolay bir iş değildir.

Yerine Koyma için salt dışsal bir gerekçelendirme sunma girişimlerinin
hepsinde ortaya çıkan bir başka sorunu ele alalım. (Bu sorun belki ilkinden
daha temeldir.) Boolos yalnızca Yerine Koyma'nın arzu edilir birçok sonucu
olduğuna dikkat çekmez. Yerine Koyma'nın ``(görünüşe göre) arzu edilmez''
hiçbir sonucu bulunmadığını da söyler. Fakat parantez içindeki ``görünüşe
göre'' kaydı kuşkusuz son derece önemlidir.

Buraya nasıl geldiğimizi hatırlayın: Naif Küme Oluşturma tutarsızlığa yol
açtı; biz de bu tutarsızlığa birikimli-yinelemeli küme anlayışını benimseyerek
yanıt verdik. Bu anlayışın yanında, umuyoruz ki tutarlılığını güvence altına
alan bir öykü bulunur. Fakat Yerine Koyma'yı bu öykünün içinden
gerekçelendiremiyorsak $\ZF$'nin tutarlı olduğuna inanmak için (henüz) hiçbir
nedenimiz yoktur. Daha doğrusu, hiç kimsenin \emph{henüz} bir çelişki
keşfetmemiş olması gibi (belki salt olumsal) bir olgu dışında, $\ZF$'nin
tutarlı olduğuna inanmak için hiçbir nedenimiz yoktur. Boolos'un yorumu tam
bu anlamda şuna indirgeniyor görünür: ``(görünüşe göre) $\ZF$ tutarlıdır.''
Tutarlılık konusunda bundan daha güçlü bir güvence istemeliyiz.

Bu sorun, Yerine Koyma'yı gerekçelendirmeye yönelik her \emph{salt} dışsal
girişimi, yani yalnızca $\ZF$'nin (bilinen) sonuçlarıyla ifade edilen her
gerekçelendirmeyi etkileyecektir. Bu nedenle Yerine Koyma için \emph{içsel}
bir gerekçelendirme, yani kümeler hakkında anlattığımız öykünün bizi bir
biçimde Yerine Koyma'ya ``zaten'' bağladığını düşündüren bir gerekçelendirme
aramak isteyeceğiz.

